% !TeX program = lualatex
% halo:begin
% {
%   "version": 1,
%   "publish": true,
%   "course": "higher-algebra-i",
%   "section": "1.4",
%   "title": "高等代数（一）1.4：矩阵与行列式展开",
%   "categories": ["study-notes"],
%   "tags": ["mathematics"],
%   "excerpt": "矩阵与行列式的区别、转置与初等变换，以及从列的多重线性推导按行（列）展开。"
% }
% halo:end
% 课堂日期：2026-09-24。1.4沿用本课程课堂笔记序号，非已核实教材节号。
% 来源：IMG_1403.HEIC、IMG_1404.HEIC；笔记 2026年9月24日.pdf及同名(2).pdf。
% 用户确认本课也讲了按行展开；单行（列）展开由上一讲迁入；例题为说明计算方法而补充。
\RequirePackage{iftex}
\RequireLuaTeX
\documentclass[UTF8, fontset=fandol, a4paper, 12pt]{ctexart}

\usepackage[margin=2.5cm]{geometry}
\usepackage{amsmath, amssymb, amsthm}
\usepackage{enumitem}
\usepackage{fancyhdr}
\usepackage[hidelinks]{hyperref}

\theoremstyle{definition}
\newtheorem{theorem}{定理}[section]
\newtheorem{proposition}[theorem]{命题}
\newtheorem{definition}[theorem]{定义}
\newtheorem{example}[theorem]{例题}
\renewcommand{\proofname}{证明}
\setlist[enumerate]{leftmargin=2em, itemsep=0.2em, topsep=0.3em}
\setlength{\headheight}{16pt}
\pagestyle{fancy}
\fancyhf{}
\fancyhead[L]{\small 高等代数（一）}
\fancyhead[R]{\small 1.4\quad 矩阵与行列式展开}
\fancyfoot[C]{\thepage}
\renewcommand{\headrulewidth}{0.3pt}
\raggedbottom

\title{矩阵与行列式展开}
\author{Yuchen Zhang}
\date{课堂日期：2026年9月24日}
\makeatletter
\renewcommand{\maketitle}{%
  \begin{center}
    {\LARGE\bfseries \@title\par}\vspace{0.6em}
    \@author\par\vspace{0.2em}
    {\small \@date\par}
  \end{center}
}
\makeatother

\begin{document}
\maketitle

\section{矩阵与行列式}

\begin{definition}
设 $P$ 为数域。把 $mn$ 个数 $a_{ij}\in P$ 排成 $m$ 行、$n$ 列的矩形数表
\[
 A=\begin{pmatrix}
 a_{11}&a_{12}&\cdots&a_{1n}\\
 a_{21}&a_{22}&\cdots&a_{2n}\\
 \vdots&\vdots&&\vdots\\
 a_{m1}&a_{m2}&\cdots&a_{mn}
 \end{pmatrix},
\]
称为一个 $m\times n$ 矩阵，简记为 $A=(a_{ij})_{m\times n}$。
其中 $a_{ij}$ 是第 $i$ 行、第 $j$ 列的元素。当 $m=n$ 时，称 $A$ 为 $n$ 阶方阵。
\end{definition}

矩阵本身是一张数表。对 $n$ 阶方阵 $A$，按排列定义计算得到的数
称为它的\textbf{行列式}，记作 $\det A$ 或 $|A|$。例如
\[
 A=\begin{pmatrix}1&2\\3&4\end{pmatrix},
 \qquad |A|=1\cdot4-2\cdot3=-2.
\]
这里 $A$ 是矩阵，$|A|$ 是数；非方阵不定义通常意义下的行列式。

将矩阵的第 $j$ 列记为
\[
 \boldsymbol\alpha_j=(a_{1j},a_{2j},\ldots,a_{mj})^{\mathsf T},
\]
则 $A=(\boldsymbol\alpha_1,\ldots,\boldsymbol\alpha_n)$。
因此，$n$ 阶行列式也可看成一个运算：输入 $n$ 个 $n$ 维列向量，输出一个数。
记号 $\det(\boldsymbol\alpha_1,\ldots,\boldsymbol\alpha_n)$ 就是以这些向量为列的方阵的行列式。

\subsection{转置}

\begin{definition}
把矩阵 $A=(a_{ij})_{m\times n}$ 的行与列互换，得到的 $n\times m$ 矩阵
称为 $A$ 的\textbf{转置矩阵}，记为 $A^{\mathsf T}$。若 $B=A^{\mathsf T}$，则 $b_{ij}=a_{ji}$。
\end{definition}

对方阵，有 $\det A^{\mathsf T}=\det A$。用排列定义可写成
\[
 \det A^{\mathsf T}
 =\sum_{\sigma\in S_n}(-1)^{\tau(\sigma)}\prod_{i=1}^n a_{\sigma(i),i}
 =\sum_{\sigma\in S_n}(-1)^{\tau(\sigma)}\prod_{k=1}^n a_{k,\sigma^{-1}(k)}.
\]
逆排列与原排列同奇偶，且 $\sigma^{-1}$ 也遍历全部排列，故末式就是 $\det A$。
这里 $S_n$ 表示所有 $n$ 元排列的集合，$\tau(\sigma)$ 表示逆序数。
因此，行与列的性质相对应。

\section{初等变换与行列式的变化}

\begin{definition}
以下三种操作称为矩阵的\textbf{初等行变换}：
\begin{enumerate}
 \item 用非零数 $c$ 乘某一行，记为 $r_i\leftarrow c r_i$；
 \item 把一行的 $c$ 倍加到另一行，记为 $r_i\leftarrow r_i+c r_j$，其中 $i\ne j$；
 \item 交换两行，记为 $r_i\leftrightarrow r_j$，其中 $i\ne j$。
\end{enumerate}
对列作同样的操作，称为\textbf{初等列变换}。
\end{definition}

这些操作对矩形矩阵也适用。若 $A$ 为方阵、变换后得到 $B$，则由行列式性质：
\[
\begin{array}{c|c}
 \text{对矩阵所作的操作}&\text{行列式的变化}\\ \hline
 r_i\leftarrow c r_i&\det B=c\det A\\
 r_i\leftarrow r_i+c r_j&\det B=\det A\\
 r_i\leftrightarrow r_j&\det B=-\det A
\end{array}
\]
列变换的结论相同。矩阵经过变换后通常改变；上表说明的是它的行列式怎样变化。

第一类变换要求 $c\ne0$，是为了能用乘 $1/c$ 的操作还原。
但作为行列式的数乘性质，$c=0$ 时等式仍成立，此时产生零行，行列式为零。
还应区分“某一行乘 $c$”与“整个矩阵乘 $c$”：后者把每一行都乘 $c$，故
\[
 \det(cA)=c^n\det A.
\]

\section{从列的多重线性到展开公式}

\subsection{固定其余列，只拆开一列}

行列式对每一列分别线性。例如固定第 $j$ 列以外的各列，有
\begin{align*}
 &\det(\boldsymbol\alpha_1,\ldots,s\boldsymbol u+t\boldsymbol v,\ldots,\boldsymbol\alpha_n)\\
 &\quad=s\det(\boldsymbol\alpha_1,\ldots,\boldsymbol u,\ldots,\boldsymbol\alpha_n)
       +t\det(\boldsymbol\alpha_1,\ldots,\boldsymbol v,\ldots,\boldsymbol\alpha_n).
\end{align*}
每次只对一列使用线性，其他列保持不动；这就是列的\textbf{多重线性}的含义。
它不能理解为对整个矩阵的线性，一般并没有 $\det(A+B)=\det A+\det B$。

设 $\boldsymbol e_i$ 为第 $i$ 个分量等于 $1$、其余分量为 $0$ 的标准单位列向量。
任意第 $j$ 列都可写成
\[
 \boldsymbol\alpha_j=\sum_{i=1}^n a_{ij}\boldsymbol e_i.
\]
代入行列式并利用第 $j$ 列的线性，得
\[
 \det A=\sum_{i=1}^n a_{ij}
 \det(\boldsymbol\alpha_1,\ldots,\boldsymbol\alpha_{j-1},
       \boldsymbol e_i,\boldsymbol\alpha_{j+1},\ldots,\boldsymbol\alpha_n).
\]
式中固定的是列号 $j$，求和指标 $i$ 遍历这一列的各行。
原行列式由此拆成 $n$ 项，每一项中的第 $j$ 列都只剩一个 $1$。

\subsection{余子式与代数余子式}

\begin{definition}
在 $n$ 阶方阵 $A$ 中（$n\ge2$），划去元素 $a_{ij}$ 所在的第 $i$ 行和第 $j$ 列，
剩下的 $n-1$ 阶方阵的行列式，称为 $a_{ij}$ 的\textbf{余子式}，记为 $M_{ij}$。
数 $A_{ij}=(-1)^{i+j}M_{ij}$ 称为它的\textbf{代数余子式}。
\end{definition}

这里的 $A_{ij}$ 是一个数，与矩阵 $A$ 本身不同。因子 $(-1)^{i+j}$ 的符号按棋盘交替排列：
\[
 \begin{pmatrix}+&-&+&\cdots\\-&+&-&\cdots\\+&-&+&\cdots\\\vdots&\vdots&\vdots&\ddots\end{pmatrix}.
\]

把上式中第 $j$ 列被替换为 $\boldsymbol e_i$ 的矩阵记为 $B$。
将 $B$ 的第 $i$ 行依次交换到第一行，再将第 $j$ 列依次交换到第一列，
共需 $(i-1)+(j-1)$ 次对换，得到
\[
 C=\begin{pmatrix}1&*\\0&B_{ij}\end{pmatrix},
 \qquad \det C=(-1)^{i+j-2}\det B.
\]
其中 $B_{ij}$ 正是从 $A$ 删去第 $i$ 行、第 $j$ 列后的矩阵，其行列顺序未变。
由排列定义，$C$ 的第一列只有第一项非零，非零乘积必须选取左上角的 $1$；
其余元素的带符号求和恰为 $\det B_{ij}=M_{ij}$。因此
\[
 \det B=(-1)^{i+j}M_{ij}=A_{ij}.
\]
所以展开式中的符号来自行、列交换的次数，而余子式来自删去那个 $1$ 所占的行与列。

\subsection{按一行或一列展开}

\begin{theorem}
设 $D=\det A$。固定第 $j$ 列，有
\[
 \boxed{D=\sum_{i=1}^n a_{ij}A_{ij}
           =\sum_{i=1}^n(-1)^{i+j}a_{ij}M_{ij}.}
\]
固定第 $i$ 行，有
\[
 \boxed{D=\sum_{j=1}^n a_{ij}A_{ij}
           =\sum_{j=1}^n(-1)^{i+j}a_{ij}M_{ij}.}
\]
\end{theorem}
第一式由上述列分解得到；第二式可对转置矩阵应用第一式得到。
这把一个 $n$ 阶行列式化为若干个 $n-1$ 阶行列式。
计算时通常选择零较多的一行或一列，以减少实际需要计算的余子式。

\newpage
\begin{example}
计算 $D=\begin{vmatrix}1&0&2\\3&4&5\\0&0&6\end{vmatrix}$。
\end{example}
\begin{proof}[解]
按第二列展开，只有 $a_{22}=4$ 非零。划去第二行、第二列，得
\[
 M_{22}=\begin{vmatrix}1&2\\0&6\end{vmatrix}=6,
 \qquad D=4(-1)^{2+2}M_{22}=24.
\]
也可以先作 $r_2\leftarrow r_2-3r_1$，得到上三角形行列式，
再取主对角线元素的乘积 $1\cdot4\cdot6=24$。
\end{proof}

\subsection{不同行（列）的对应乘积之和}

\begin{proposition}
某一行的元素与另一行对应位置的代数余子式相乘，再求和，结果为零。
连同展开公式，有
\[
 \sum_{r=1}^n a_{hr}A_{ir}
 =\begin{cases}D,&h=i,\\0,&h\ne i.\end{cases}
\]
相应的列形式为
\[
 \sum_{r=1}^n a_{rk}A_{rj}
 =\begin{cases}D,&k=j,\\0,&k\ne j.\end{cases}
\]
\end{proposition}
\begin{proof}
当 $h=i$ 时，就是按第 $i$ 行展开。
当 $h\ne i$ 时，把 $A$ 的第 $i$ 行替换为第 $h$ 行，得到两行相同的矩阵，行列式为零。
沿替换后的第 $i$ 行展开；由于计算余子式时会删去这一行，代数余子式仍为原来的 $A_{ir}$，
故展开式正是 $\sum_r a_{hr}A_{ir}=0$。列的结论同理。
\end{proof}

\end{document}
